%% 15장 알고리즘 분석, 정렬과 선택
\chapter{알고리즘 분석, 정렬과 선택}
\label{ch:sorting}

\abstract{정렬은 전산학에서 가장 많이 연구된 문제이고, 그 연구에서 알고리즘을 분석하는 방법 대부분이 나왔다. 이 장에서는 정렬과 선택을 통해 그 방법들을 배운다. 먼저 비교로 정렬하는 모든 알고리즘이 넘을 수 없는 하한 $\log_2n!$을 결정 나무로 증명하고, 병합 정렬과 퀵정렬이 그 하한에 얼마나 가까운지를 각각 최악의 경우와 평균의 경우로 분석한다. 퀵정렬의 약점을 힙 정렬로 막는 인트로 정렬을 본 뒤, 비교를 하지 않아 하한을 피해 가는 계수 정렬과 기수 정렬로 넘어간다. 이어서 분기 없이 정해진 비교만으로 정렬하는 정렬 네트워크와 그 정확성을 다루는 0-1 원리를 증명하고, 마지막으로 $k$번째로 작은 값을 찾는 선택 문제를 평균과 최악의 경우 모두 선형 시간에 푸는 방법과 이분 탐색을 다룬다. stardust의 배열 정렬은 이 알고리즘들을 레인의 길이와 자료형에 따라 골라 쓴다.}

\section{알고리즘의 분석}
\label{sec:analysis}

알고리즘의 비용은 입력의 크기 $n$의 함수로 나타낸다. 무엇을 셀지는 문제에 따라 정하는데, 정렬에서는 원소끼리의 비교 횟수가 자연스러운 척도이다. 비교가 가장 비싼 연산이라는 보장은 없지만, 비교 정렬에서 다른 연산의 횟수는 대개 비교 횟수에 비례하기 때문이다. 같은 크기의 입력이라도 비용은 입력에 따라 다르므로, 가장 나쁜 입력일 때의 \term{최악의 경우}{worst case} 비용과, 입력에 확률분포를 주었을 때의 \term{평균의 경우}{average case} 비용을 따로 분석한다. 무작위성을 쓰는 알고리즘에서는 입력을 고정하고 알고리즘의 무작위 선택에 대한 기댓값을 따지는데, 이것은 어떤 입력에도 성립하므로 평균의 경우보다 강한 보장이다. 결과는 3장의 점근 기호로 적는다.

13장에서 보았듯이 실제 실행 시간은 연산의 수만으로 정해지지 않는다. 그래도 점근 분석은 알고리즘이 큰 입력에서 어떻게 행동하는지를 말해 주는 가장 기본적인 도구이고, 상수 인자와 메모리 접근 패턴은 그 위에서 따로 따진다. 이 장에서도 비교 횟수의 분석과 함께 실제로 어느 알고리즘이 빠른지를 본다.

\section{비교 정렬의 하한}
\label{sec:lower-bound}

\begin{theorem}
\label{thm:sort-lower-bound}
서로 다른 $n$개의 원소를 비교만으로 정렬하는 알고리즘은 최악의 경우 적어도 $\lceil\log_2n!\rceil$번 비교한다. 입력의 $n!$가지 순서가 모두 같은 확률이면 평균 비교 횟수도 적어도 $\log_2n!$이다.
\end{theorem}

\begin{proof}
알고리즘의 동작을 \term{결정 나무}{decision tree}로 나타내자. 안쪽 마디는 비교 하나이고, 비교 결과에 따라 왼쪽이나 오른쪽 자식으로 가며, 잎은 알고리즘이 내놓는 순열이다. 서로 다른 두 입력 순서는 서로 다른 순열로 정렬되어야 하므로 잎은 적어도 $n!$개이다. 높이가 $h$인 이진 나무의 잎은 많아야 $2^h$개이므로 $2^h\ge n!$, 곧 $h\ge\log_2n!$이다. 평균에 대한 주장은 잎이 $L$개인 이진 나무에서 잎의 깊이의 평균이 $\log_2L$ 이상이라는 사실에서 나온다. 이 사실은 Kraft 부등식 $\sum_{\text{잎}}2^{-d}\le1$과 Jensen 부등식에서 $2^{-\bar d}\le(1/L)\sum2^{-d}\le1/L$이 성립한다는 것으로 보일 수 있다.
\end{proof}

3장의 Stirling 공식에서 $\log_2n!=n\log_2n-n\log_2e+O(\log n)$이므로 하한은 대략 $n\log_2n-1.443n$이다. 구체적인 값을 보면 $n=8$에서 $\lceil\log_28!\rceil=16$, $n=16$에서 $45$, $n=10^6$에서 약 $1.85\times10^7$이다. 이 하한은 비교만 쓴다는 가정에서 나온 것이므로, 원소의 값 자체를 쓰는 기수 정렬은 이 하한의 적용을 받지 않는다(\cref{sec:radix}).

\section{비교 정렬}
\label{sec:comparison-sorts}

\subsection{삽입 정렬과 병합 정렬}

\term{삽입 정렬}{insertion sort}은 앞쪽의 이미 정렬된 부분에 원소를 하나씩 제자리를 찾아 넣는다. 최악의 경우 $n(n-1)/2$번 비교하지만, 반복이 단순하고 메모리를 차례로 접근하며 거의 정렬된 입력에서는 선형 시간이므로, 원소가 십여 개 이하인 짧은 구간에서는 어떤 정교한 알고리즘보다도 빠르다. 그래서 거의 모든 실용적인 정렬 알고리즘은 재귀를 짧은 구간에서 멈추고 삽입 정렬로 마무리한다.

\term{병합 정렬}{merge sort}은 배열을 반으로 나누어 각각을 재귀적으로 정렬한 뒤 두 정렬된 열을 하나로 병합한다. 길이 $a$와 $b$인 두 열의 병합은 앞에서부터 작은 쪽을 꺼내는 일이므로 많아야 $a+b-1$번 비교한다. 따라서 비교 횟수 $C(n)$은 $C(n)\le C(\lceil n/2\rceil)+C(\lfloor n/2\rfloor)+n-1$을 만족하고, 귀납법으로 $C(n)\le n\lceil\log_2n\rceil$을 얻는다. 하한 $n\log_2n-1.443n$과 비교하면 최악의 경우에도 거의 최적이다.

병합 정렬의 또 다른 장점은 \term{안정성}{stability}이다. 병합할 때 두 열의 원소가 같으면 왼쪽 열의 것을 먼저 꺼내면, 원래 입력에서 앞에 있던 원소가 결과에서도 앞에 온다. 1장에서 보았듯이 안정 정렬은 여러 키로 정렬하거나 순위를 매길 때 필수적이다. 약점은 병합에 원래 배열만큼의 추가 메모리가 필요하다는 것이다.

\subsection{퀵정렬}

Hoare의 \term{퀵정렬}{quicksort}은 원소 하나를 \term{피벗}{pivot}으로 골라, 피벗보다 작은 원소들을 왼쪽에, 큰 원소들을 오른쪽에 모으는 \term{분할}{partition}을 한 뒤 양쪽을 재귀적으로 정렬한다\cite{hoare1962}. 분할은 배열을 제자리에서 바꾸므로 추가 메모리가 거의 필요 없고, 메모리를 양 끝에서 차례로 훑으므로 캐시에 잘 맞는다. 피벗이 늘 가장 작은 원소이면 분할이 $1$과 $n-1$로 갈려 비교가 $\Theta(n^2)$번이지만, 평균적으로는 그렇지 않다.

\begin{theorem}
\label{thm:quicksort-average}
서로 다른 $n$개의 원소에서 피벗을 매번 균등한 무작위로 고르는 퀵정렬의 비교 횟수의 기댓값은
\[
2(n+1)H_n-4n,\qquad H_n=\sum_{k=1}^n\frac1k
\]
이고, 이것은 $2n\ln n+O(n)\approx1.386\,n\log_2n$이다.
\end{theorem}

\begin{proof}
원소들을 정렬한 순서로 $z_1<z_2<\cdots<z_n$이라 하자. 두 원소는 둘 중 하나가 피벗일 때만 비교되고, 같은 두 원소가 두 번 비교되는 일은 없다. $i<j$에 대해 $z_i$와 $z_j$가 비교되는지를 보자. 집합 $Z_{ij}=\{z_i,\dots,z_j\}$의 원소들은 그 가운데 하나가 피벗으로 뽑히기 전까지는 늘 같은 쪽으로 분할되므로 같은 부분 배열에 머문다. $Z_{ij}$에서 처음 뽑히는 피벗이 $z_i$나 $z_j$이면 둘은 비교되고, 그 사이의 다른 원소이면 둘은 서로 다른 쪽으로 갈려 영영 비교되지 않는다. 피벗은 균등하게 뽑으므로 $Z_{ij}$에서 처음 뽑히는 원소는 $Z_{ij}$의 $j-i+1$개 원소 가운데 균등하게 하나이다. 따라서 비교될 확률은 $2/(j-i+1)$이고, 기댓값의 선형성에서
\[
\sum_{i<j}\frac2{j-i+1}=2\sum_{k=2}^n\frac{n-k+1}k=2\Bigl((n+1)(H_n-1)-(n-1)\Bigr)=2(n+1)H_n-4n
\]
이다. 둘째 등호는 차이가 $k-1$인 쌍 $(i,j)$가 $n-k+1$개라는 데서, 셋째 등호는 $\sum_{k=2}^n(n+1)/k=(n+1)(H_n-1)$에서 나온다. $H_n=\ln n+O(1)$이므로 마지막 주장도 성립한다.
\end{proof}

평균 비교 횟수는 하한보다 약 $39\%$ 많다. $n=10^6$이면 하한이 $1.85\times10^7$이고 퀵정렬의 기댓값은 $2.48\times10^7$이다. 그런데도 퀵정렬이 병합 정렬보다 빠른 경우가 많은 것은 안쪽 반복이 단순하고 데이터를 옮기는 횟수가 적기 때문이다. 실제 구현은 피벗을 무작위로 고르는 대신 처음, 가운데, 끝 원소의 중앙값을 고르는데, 이 방법은 이미 정렬된 입력에서도 좋은 피벗을 주고 평균 비교 횟수의 상수도 줄인다\cite{sedgewick1977quicksort}.

\subsection{인트로 정렬}

퀵정렬의 최악의 경우 $\Theta(n^2)$은 이론적인 걱정만이 아니다. 피벗을 고르는 규칙을 아는 사람은 그 규칙을 무너뜨리는 입력을 만들 수 있고, 실제로 그런 입력으로 서버를 느리게 만드는 공격이 있었다. Musser의 \term{인트로 정렬}{introsort}은 퀵정렬의 재귀 깊이를 세다가 $2\log_2n$을 넘으면 그 부분 배열을 \term{힙 정렬}{heapsort}로 정렬한다\cite{musser1997}. 힙 정렬은 최악의 경우에도 $O(n\log n)$이므로(16장에서 힙을 다룬다\cite{williams1964heapsort}), 인트로 정렬은 평균적으로는 퀵정렬만큼 빠르고 최악의 경우에도 $O(n\log n)$을 보장한다. C++ 표준 라이브러리의 \code{std::sort}가 대개 이 방법을 쓰며, 최근에는 흔한 입력 패턴을 알아채 더 빨라지는 pattern-defeating 퀵정렬 같은 변형도 쓰인다\cite{peters2021pdq}.

\section{비교하지 않는 정렬}
\label{sec:radix}

\subsection{계수 정렬}

원소가 $0$부터 $K-1$까지의 정수이면 비교 없이 정렬할 수 있다. 먼저 각 값이 몇 번 나오는지 세고, 그 개수를 누적해 각 값이 결과에서 시작할 위치를 구한 뒤, 입력을 처음부터 훑으며 각 원소를 자기 값의 다음 자리에 놓는다. 이 \term{계수 정렬}{counting sort}은 $O(n+K)$ 시간이 걸리고 입력을 앞에서부터 훑으며 같은 값을 차례로 놓으므로 안정적이다.

\subsection{기수 정렬}

$K$가 크면 계수 정렬의 표가 너무 크다. 64비트 키라면 $K=2^{64}$이다. \term{기수 정렬}{radix sort}은 키를 $8$비트씩 자리로 나누어 계수 정렬을 자리마다 되풀이한다. 가장 낮은 자리부터 시작하는 LSD(least significant digit) 기수 정렬이 올바른 이유는 다음과 같다.

\begin{theorem}
\label{thm:lsd-radix}
키를 $d$개의 자리 $k=(a_{d-1},\dots,a_1,a_0)$로 나누어, 가장 낮은 자리 $a_0$부터 $a_{d-1}$까지 각 자리를 키로 하는 안정 정렬을 차례로 하면, 결과는 키 전체의 사전식 순서, 곧 키의 수치 순서로 정렬되고 같은 키끼리는 입력 순서가 유지된다.
\end{theorem}

\begin{proof}
$t$번째 패스가 끝난 뒤 배열이 아래쪽 $t$개 자리 $(a_{t-1},\dots,a_0)$의 사전식 순서로 정렬되어 있고 그 순서가 같은 원소끼리는 입력 순서라는 것을 $t$에 대한 귀납법으로 보인다. $t=0$에서는 자명하다. $t+1$번째 패스는 자리 $a_t$로 안정 정렬한다. 두 원소의 $a_t$가 다르면 그 자리가 순서를 정하고, 이것은 아래쪽 $t+1$개 자리의 사전식 순서와 일치한다. $a_t$가 같으면 안정성에 의해 패스 전의 순서가 유지되는데, 귀납 가정에서 그 순서는 아래쪽 $t$개 자리의 사전식 순서이고 그것마저 같으면 입력 순서이다. 따라서 $t+1$에서도 주장이 성립한다.
\end{proof}

자리 하나가 $8$비트이면 표는 $256$칸이고, $64$비트 키는 패스 $8$번으로 정렬된다. 비교 정렬이 $n\log_2n$에 비례하는 반면 기수 정렬은 $8n$ 정도의 일을 하므로 $n$이 크면 훨씬 빠르다. 부호 있는 정수와 부동소수점 수는 1장의 키 사상으로 순서가 같은 부호 없는 정수로 바꾸어 정렬한다. 한 가지 기교를 더하면 더 빨라지는데, 처음에 모든 키를 한 번 훑어 모든 키에서 같은 값을 갖는 자리를 찾아 그 패스를 건너뛰는 것이다. 작은 범위의 정수라면 위쪽 자리가 모두 같으므로 패스 대부분이 사라진다.

\begin{instardust}[\code{docs/design/array.md} §8.4]
stardust는 길이 $1024$ 이상의 int64나 float64 값, 그리고 길이 $128$ 이상의 (값, 인덱스) 쌍을 $8$비트 자리의 LSD 기수 정렬로 정렬한다. 키는 1장의 사상 그대로이고, NaN은 미리 떼어 놓으며, $-0.0$은 키를 만들 때 $+0.0$으로 바꾸어 두 영이 동점이 되게 한다. 최솟값과 최댓값을 먼저 구해 그 차이의 자리만 정렬하는데, 부호가 섞인 정수는 부호 비트를 뒤집은 키의 위쪽 바이트가 오락가락해서 건너뛸 자리가 없으므로 최솟값을 빼서 범위를 줄인다. 설계 문서에 따르면 $[-1000,1000)$ 범위의 정수 백만 개에서 이렇게 하면 두 자리로 끝나 $10.5$ms가 $2.4$ms로 줄었다. 기수 정렬은 안정적이므로 \code{stable=False}여도 같은 알고리즘을 쓴다. 안정 정렬의 결과는 불안정 정렬의 올바른 결과이기도 하기 때문이다.
\end{instardust}

\section{정렬 네트워크}
\label{sec:sorting-networks}

\subsection{비교기의 나열}

\term{비교기}{comparator} $[i:j]$($i<j$)는 두 자리 $i$와 $j$의 값을 비교해 작은 것을 $i$에, 큰 것을 $j$에 둔다. 비교기들을 정해진 순서로 나열한 것을 \term{비교 네트워크}{comparator network}라 하고, 모든 입력을 정렬하는 비교 네트워크를 \term{정렬 네트워크}{sorting network}라 한다. 정렬 네트워크는 입력에 상관없이 같은 비교를 같은 순서로 하므로 분기가 없고, 비교기 하나는 최솟값과 최댓값 두 연산이어서 13장에서 본 SIMD로 여러 열을 한꺼번에 정렬할 수 있다. 짧은 배열에서는 이것이 큰 장점이다.

비교기의 수는 비교 정렬의 하한 이상이어야 한다. 입력 $8$개의 정렬 네트워크는 비교기가 적어도 $19$개 필요하고 $19$개로 충분하다는 것이 알려져 있으며\cite{knuth1998taocp3}, Batcher의 odd-even 병합 네트워크는 입력 $2^k$개를 비교기 $(k^2-k+4)2^{k-2}-1$개로 정렬해 입력 $8$개에서 $19$개, $16$개에서 $63$개를 쓴다\cite{batcher1968}.

\subsection{0-1 원리}

어떤 비교 네트워크가 정렬 네트워크인지 확인하려면 원칙적으로 모든 입력을 시험해야 할 것 같지만, 다음 정리 덕분에 $0$과 $1$만으로 된 입력 $2^n$가지만 시험하면 된다.

\begin{theorem}[0-1 원리]
\label{thm:zero-one}
$n$개의 입력을 받는 비교 네트워크가 $0$과 $1$로 된 모든 입력을 정렬하면, 그 네트워크는 전순서 집합의 원소로 된 모든 입력을 정렬한다.
\end{theorem}

\begin{proof}
먼저 단조 증가함수 $f$에 대해 네트워크와 $f$가 교환된다는 것, 곧 입력의 각 원소에 $f$를 적용한 뒤 네트워크를 통과시킨 결과가 네트워크를 통과시킨 뒤 $f$를 적용한 결과와 같다는 것을 보인다. 비교기 하나에 대해 $f$가 단조이면 $f(\min(a,b))=\min(f(a),f(b))$이고 $\max$도 마찬가지이므로 성립하고, 네트워크는 비교기의 합성이므로 전체에서도 성립한다. 이제 네트워크가 어떤 입력 $x$를 정렬하지 못한다고 하자. 그러면 출력의 어떤 자리 $i<j$에서 $y_i>y_j$이다. $f(t)=[t\ge y_i]$로 두면 $f$는 단조 증가이고 값이 $0$ 또는 $1$이다. 입력 $f(x)$는 $0$과 $1$로 된 입력이고, 교환성에서 그 출력은 $f(y)$인데, $f(y_i)=1$이고 $f(y_j)=0$이므로 $f(y)$는 정렬되어 있지 않다. 이것은 가정에 모순이다.
\end{proof}

\begin{instardust}[\code{docs/design/array.md} §8.2, §8.3]
stardust는 길이 $16$ 이하의 레인을 Batcher의 $16$입력 odd-even 병합 네트워크로 정렬한다. 레인의 길이 $m$이 $16$보다 작으면 $m$ 이상의 자리를 건드리는 비교기를 뺀 네트워크를 쓰는데, 빈 자리를 $+\infty$로 채운 것과 같으므로 여전히 올바른 정렬 네트워크이다. $m$마다 비교기 목록을 컴파일할 때 만들어 펼쳐 두어 값이 레지스터에 머물고, 비교기 하나는 비교 하나와 선택 둘이어서 분기가 없다. 길이 $32$에서 $1023$ 사이의 값 레인에서는 NEON 레지스터 여덟 개에 원소 $16$개를 싣고 $8$입력 최적 네트워크(비교기 $19$개)를 세로로 적용해 두 열을 동시에 정렬하는데, 설계 문서는 이 네트워크의 정확성을 \cref{thm:zero-one}에 따라 $0$과 $1$로 된 입력 $256$가지로 확인했다고 적는다. 길이가 $17$에서 $31$ 사이이거나 자료형이 다른 짧은 레인은 병합 정렬(안정)이나 인트로 정렬(불안정)로 간다.
\end{instardust}

\section{선택}
\label{sec:selection}

\subsection{빠른 선택}

$n$개의 원소 가운데 $k$번째로 작은 것, 예컨대 중앙값을 찾는 \term{선택}{selection} 문제는 정렬하면 풀리지만 정렬보다 적은 일로 풀 수 있다. Hoare의 \term{빠른 선택}{quickselect}은 퀵정렬처럼 분할하되 찾는 원소가 있는 쪽 하나만 재귀한다\cite{hoare1961find}. 피벗을 균등하게 고르면 비교 횟수의 기댓값은 $O(n)$인데, 간단한 상한은 다음과 같이 얻는다. 분할 하나는 $n-1$번 비교하고, 피벗의 순위가 가운데 절반, 곧 $[n/4,3n/4]$에 들어갈 확률은 약 $1/2$이며 그때 남는 부분은 $3n/4$ 이하이다. 따라서 남는 크기가 $3/4$배 이하로 줄어들기까지 평균 두 번의 분할이면 되고, 기댓값은 많아야 $2n(1+3/4+(3/4)^2+\cdots)=8n$ 정도이다. 정밀하게 계산하면 중앙값을 찾는 데 드는 비교의 기댓값은 $2(1+\ln2)n\approx3.39n$이다\cite{knuth1998taocp3}. Floyd와 Rivest는 표본에서 피벗을 두 개 골라 찾는 원소를 좁은 구간으로 몰아넣어 이것을 $n+\min(k,n-k)+o(n)$까지 줄였다\cite{floyd1975}.

\subsection{최악의 경우 선형 시간}

빠른 선택도 운이 나쁘면 $\Theta(n^2)$이다. Blum, Floyd, Pratt, Rivest, Tarjan은 피벗을 다음처럼 고르면 최악의 경우에도 선형 시간임을 보였다\cite{blum1973}. 원소들을 다섯 개씩 묶어 각 묶음의 중앙값을 구하고, 그 $n/5$개 중앙값들의 중앙값을 재귀적으로 구해 피벗으로 쓴다. 이 피벗보다 작은 원소는 적어도 묶음 중앙값의 절반, 곧 $n/10$개의 묶음에서 각각 세 개씩이므로 $3n/10$개 이상이고, 큰 원소도 마찬가지이다. 따라서 분할 뒤 남는 쪽은 $7n/10$ 이하이고, 비용은
\[
T(n)\le T(n/5)+T(7n/10)+cn
\]
을 만족한다. $1/5+7/10<1$이므로 $T(n)\le10cn$임을 귀납법으로 확인할 수 있다. 상수가 커서 실제로는 잘 쓰지 않지만, 빠른 선택이 깊이 들어가면 이 방법으로 바꾸는 인트로 선택이 최악의 경우를 막는 데 쓰인다\cite{musser1997}.

\subsection{여러 순서 통계량}

분위수 여러 개를 한꺼번에 구할 때는 순서 통계량 여러 개가 필요하다. 원하는 순위 $k_1<\cdots<k_m$을 왼쪽부터 하나씩 선택하면 $O(mn)$이지만, 가운데 순위를 먼저 선택한 뒤 그보다 작은 순위는 왼쪽 부분에서, 큰 순위는 오른쪽 부분에서만 재귀적으로 선택하면, 재귀의 각 층에서 하는 일이 $O(n)$이고 층이 $\log_2m$개이므로 $O(n\log m)$이다.

\begin{instardust}[\code{csrc/array/ops/quantile.cpp}]
배열의 \code{quantile}과 \code{median}은 레인마다 값을 작업 버퍼에 복사해 제자리에서 선택한다. 필요한 순서 통계량들은 위의 방법대로 가운데 것부터 선택하고, 범위의 양 끝에 놓인 순위는 선택 대신 최솟값이나 최댓값을 한 번 훑어 구한다. 주석에는 측정에서 정한 문턱도 적혀 있다. 원소 하나당 비용으로 NEON 정렬은 길이 $16$, $256$, $1000$, $4096$에서 $2.4$, $4.1$, $5.6$, $7.3$ns였고 \code{nth\_element}는 $8.3$, $8.0$, $7.8$, $7.5$ns였으므로, 짧은 레인은 아예 정렬한다. 모든 비교는 1장에서 본 $-0.0<+0.0$의 전순서로 하므로, 정렬하든 선택하든 같은 값이 부호까지 같게 나온다.
\end{instardust}

\section{이분 탐색}
\label{sec:binary-search}

정렬된 배열 $a_0\le a_1\le\cdots\le a_{n-1}$에서 값 $v$의 자리를 찾는 \term{이분 탐색}{binary search}은 가장 단순한 알고리즘이지만 올바르게 쓰기가 놀랄 만큼 어렵다. 경계를 하나씩 어긋나게 쓰는 실수가 흔하기 때문이다. 이런 실수를 피하는 가장 좋은 방법은 \textbf{불변식}을 정하고 지키는 것이다. ``$a_i\le v$인 $i$의 개수''를 구한다고 하자. 구간 $[lo,hi)$를 두고 언제나 $a_i\le v$($i<lo$)이고 $a_i>v$($i\ge hi$)라는 불변식을 유지하면, 처음에 $lo=0$, $hi=n$이고, 가운데 $m$에서 $a_m\le v$이면 $lo\gets m+1$, 아니면 $hi\gets m$으로 둔다. 두 경우 모두 불변식이 유지되고 구간이 줄어들므로, $lo=hi$가 되면 그 값이 답이다. 비교 횟수는 $\lceil\log_2(n+1)\rceil$ 이하이며, 이것은 결정 나무의 논법으로 보면 최적이다.

$a_i<v$인 개수와 $a_i\le v$인 개수는 같은 값이 여러 번 나올 때 다르다. numpy의 \code{searchsorted}는 앞의 것을 \code{side='left'}, 뒤의 것을 \code{side='right'}로 구별한다. 찾는 값들이 정렬되어 있으면 직전 값의 결과에서 시작해 거리를 $1,2,4,\dots$로 늘려 가며 경계를 넘는 곳을 찾은 뒤 마지막 구간에서만 이분 탐색하는 \term{지수 탐색}{exponential search}이 더 빠르다. 13장에서 본 분기 없는 이분 탐색과 함께, stardust의 구간 탐색이 이 두 방법을 쓴다.

\section{맺음}

정렬은 하나의 문제에 대해 이론적 하한, 평균과 최악의 분석, 무작위화, 비교를 넘어서는 방법, 하드웨어에 맞춘 구조가 모두 등장하는 드문 예이다. 실제 라이브러리는 이 가운데 하나를 고르지 않고 입력의 크기와 자료형에 따라 여러 알고리즘을 섞어 쓴다. 짧은 구간에는 정렬 네트워크와 삽입 정렬을, 긴 정수와 실수에는 기수 정렬을, 그 밖에는 병합 정렬이나 인트로 정렬을 쓰며, 그 경계는 측정으로 정한다. 다음 장에서는 힙 정렬의 바탕인 힙을 비롯해 이런 알고리즘들이 쓰는 자료구조를 다룬다.

정렬과 탐색의 고전적인 참고서는 Knuth의 \cite{knuth1998taocp3}이고, 알고리즘 분석 전반은 \cite{cormen2022clrs}가 표준적인 교과서이다.

\chapterreferences
